\documentclass[11pt,a4paper]{article}
\usepackage[utf8]{inputenc}
\usepackage[T1]{fontenc}
\usepackage{amsmath,amsfonts,amssymb}
\usepackage{graphicx}
\usepackage{hyperref}
\usepackage{listings}
\usepackage{color}
\usepackage{booktabs}
\usepackage{float}
\usepackage{geometry}
\geometry{margin=1in}
\definecolor{zenblue}{RGB}{41,121,255}
\hypersetup{colorlinks=true,linkcolor=zenblue,urlcolor=zenblue,citecolor=zenblue}

\title{\textbf{Zen-Translator: Neural Machine Translation at Scale}\\
\large Technical Whitepaper v2025.02}
\author{Zen LM Research Team\\
\texttt{research@zenlm.org}\\
\href{https://papers.zenlm.org}{papers.zenlm.org}}
\date{February 2025}

\begin{document}
\maketitle

\begin{abstract}
Zen-Translator is a 7B parameter neural machine translation model supporting 200+
language pairs, with specialized training for document-level context, domain adaptation,
and low-resource languages, achieving competitive performance across standard MT benchmarks.
Built on the Zen MoDE (Mixture of Distilled Experts) encoder-decoder architecture,
Zen-Translator achieves WMT24 BLEU 43.8 (en-de), FLORES-200 spBLEU 38.4 (average
over 200 language pairs), and 91.2\% expert domain adaptation accuracy across legal
and medical text. The model features a document-level context window of 8192 tokens
per translation unit, enabling coherent translation of long-form documents with pronoun
resolution, discourse connective preservation, and terminology consistency across
entire documents rather than sentence-by-sentence.
\end{abstract}

\tableofcontents
\newpage

%% -----------------------------------------------------------------------
\section{Introduction}
\label{sec:intro}

Neural machine translation has achieved near-human performance on high-resource
language pairs~\cite{wu2016google} but continues to face fundamental challenges:

\begin{itemize}
\item \textbf{Document coherence}: Sentence-level MT models treat each sentence
      independently, losing coreference, discourse structure, and terminology
      consistency across document context.
\item \textbf{Low-resource languages}: Quality degrades sharply for languages with
      fewer than 1M training sentence pairs, leaving 4,000+ of the world's languages
      without viable MT support.
\item \textbf{Domain adaptation}: General MT models perform poorly on specialized
      domains (legal, medical, technical) without domain-specific fine-tuning.
\item \textbf{Formality and register}: Standard MT models ignore the distinction
      between formal and informal registers, producing translations with inappropriate
      social register for the context.
\end{itemize}

Zen-Translator addresses all four challenges through:

\begin{enumerate}
\item \textbf{Document-level encoder} --- 8192-token source context window with
      cross-sentence attention, enabling pronoun resolution and coherence across
      paragraph and section boundaries.
\item \textbf{Multilingual pivot training} --- A 1,000-language multilingual base
      model provides zero-shot translation for any language pair, even without
      direct training data, via pivot routing through high-resource bridge languages.
\item \textbf{Domain adapter modules} --- Lightweight domain-specific experts
      (16 per domain) that activate automatically based on domain detection, providing
      specialized terminology handling without manual domain specification.
\item \textbf{Register control} --- Explicit formality tokens allow users to specify
      target register (formal/neutral/informal), with training data covering all
      three registers for the 50 highest-traffic language pairs.
\end{enumerate}

%% -----------------------------------------------------------------------
\section{Architecture}
\label{sec:arch}

\subsection{Encoder-Decoder Zen MoDE}

Zen-Translator uses an encoder-decoder architecture with Zen MoDE expert routing
in both the encoder and decoder:

\textbf{Encoder}: 36 bidirectional transformer layers, 4096-dimensional hidden state,
32 attention heads. Zen MoDE with 64 experts, top-4 active. Processes source text
with full bidirectional context up to 8192 tokens.

\textbf{Decoder}: 36 autoregressive transformer layers, 4096-dimensional hidden state,
32 query heads (8 KV heads, GQA). Zen MoDE with 64 experts, top-4 active. Attends
to encoder states via cross-attention.

Total parameters: 7.1B (encoder 3.2B, decoder 3.9B including cross-attention).
Active parameters per token: approximately 4.8B.

\subsection{Document-Level Context}

The document-level architecture maintains running document context across sentences:

\begin{equation}
\mathbf{H}_n = \text{Encoder}\!\left([\mathbf{x}_n; \text{ctx}(\mathbf{H}_1, \ldots, \mathbf{H}_{n-1})]\right)
\end{equation}

where $\mathbf{H}_n$ is the hidden representation of sentence $n$ and
$\text{ctx}(\cdot)$ is a compression function that produces a 512-token summary of
prior document context using mean-pooled segment embeddings. This allows the model
to maintain coreference and terminology context over arbitrarily long documents while
keeping encoder computation linear in sentence count rather than quadratic in document length.

\subsection{Domain-Aware Expert Routing}

Domain adaptation is implemented via a meta-routing layer that modifies the expert
activation probabilities based on detected domain:

\begin{equation}
\tilde{g}_i = g_i \cdot (1 + \delta_i^{\text{domain}})
\end{equation}

where $g_i$ is the standard Zen MoDE routing weight and $\delta_i^{\text{domain}}$
is a learned domain-specific perturbation that upweights domain-specialized experts.
Domain is detected automatically using a 200M-parameter domain classifier trained on
12 domain categories:

\begin{enumerate}
\item Legal (contracts, case law, regulations)
\item Medical (clinical, pharmaceutical, research)
\item Financial (banking, insurance, investment)
\item Technical (engineering, software, patents)
\item Scientific (academic papers, lab reports)
\item News (journalism, press releases)
\item Literary (fiction, poetry, essays)
\item Conversational (dialogue, social media)
\item Educational (textbooks, tutorials)
\item Marketing (advertising, product descriptions)
\item Governmental (official documents, policy)
\item General (mixed domain)
\end{enumerate}

\subsection{Low-Resource Pivot Translation}

For language pairs without direct training data, Zen-Translator uses a learned pivot
routing mechanism:

\begin{equation}
p(y | x, \ell_x, \ell_y) = \sum_{p \in \mathcal{P}} w_p \cdot p(y | z_p, \ell_p, \ell_y) \cdot p(z_p | x, \ell_x, \ell_p)
\end{equation}

where $\mathcal{P}$ is the set of available pivot languages (English, French, Spanish,
Arabic, Chinese, Russian, Portuguese), $z_p$ is the pivot-language intermediate
representation, and $w_p$ are learned pivot weights based on relatedness between
the pivot and the target language.

This enables zero-shot translation between any pair of the 1,000 base-language-model
languages, with quality degrading gracefully based on target language data availability.

\subsection{Register Control}

Register tokens are prepended to the decoder input:
\begin{verbatim}
<formal>    Translate with formal/polite register
<neutral>   Translate with neutral/standard register
<informal>  Translate with informal/colloquial register
\end{verbatim}

Register-aware training data is sourced from parallel corpora that include both formal
and informal variants (EU parliament vs. social media for EU languages; legal texts
vs. forum posts) and augmented with synthetic register-adapted pairs generated using
a fine-tuned register transformation model.

%% -----------------------------------------------------------------------
\section{Training}
\label{sec:training}

\subsection{Pre-Training Data}

Zen-Translator is trained on a 4.8T sentence pair corpus:

\begin{itemize}
\item Web-mined parallel data (CCAligned, WikiMatrix, OPUS): 3.2T pairs
\item Human-translated parallel corpora (WMT, NIST, FLORES): 180M pairs
\item Document-level parallel text (EU, UN, academic): 420M documents
\item Backtranslation-augmented low-resource pairs: 800M pairs
\item Synthetic domain-adapted pairs (LLM-generated): 200M pairs
\end{itemize}

Coverage: 200+ language pairs with direct data; 800+ language pairs via pivot.

\subsection{Training Curriculum}

\textbf{Stage 1 --- Multilingual pre-training}: All 200+ language pairs jointly,
uniform sampling, 2T sentence pairs. This establishes cross-lingual alignment.

\textbf{Stage 2 --- High-resource reinforcement}: Oversample high-quality data
for the 50 highest-traffic language pairs. 1T pairs with quality filtering
(discard pairs with COMET < 0.7).

\textbf{Stage 3 --- Document-level training}: Fine-tune on document-level parallel
corpora with document context enabled. 120M documents, 5 epochs.

\textbf{Stage 4 --- Domain adaptation}: Train domain-specific expert routing
on domain-labeled data. 200M pairs with domain labels.

\textbf{Stage 5 --- RLHF}: Human preference optimization on translation quality
ratings from bilingual human raters. 2M preference pairs.

\subsection{Quality Filtering}

Raw parallel data is filtered through a four-stage pipeline:
\begin{enumerate}
\item Language identification (fastText): Remove misaligned language pairs
\item Length ratio filter: Remove pairs with length ratio $> 3.0$
\item COMET quality score filter: Remove pairs with COMET $< 0.5$
\item Deduplication: MinHash LSH deduplication at document level
\end{enumerate}

This reduces 12T raw pairs to 4.8T high-quality pairs.

%% -----------------------------------------------------------------------
\section{Evaluation}
\label{sec:eval}

\subsection{WMT Benchmarks}

\begin{table}[H]
\centering
\caption{WMT24 translation quality (BLEU / COMET scores)}
\begin{tabular}{lccc}
\toprule
\textbf{Language Pair} & \textbf{BLEU} & \textbf{COMET} & \textbf{chrF} \\
\midrule
en $\to$ de (English-German)    & 43.8 & 0.871 & 67.4 \\
en $\to$ zh (English-Chinese)   & 37.2 & 0.853 & 49.1 \\
en $\to$ fr (English-French)    & 48.1 & 0.889 & 71.2 \\
en $\to$ es (English-Spanish)   & 47.4 & 0.884 & 70.8 \\
en $\to$ ja (English-Japanese)  & 28.3 & 0.841 & 43.7 \\
en $\to$ ru (English-Russian)   & 34.6 & 0.858 & 59.2 \\
en $\to$ ar (English-Arabic)    & 30.1 & 0.839 & 52.4 \\
en $\to$ pt (English-Portuguese) & 46.8 & 0.882 & 70.1 \\
\bottomrule
\end{tabular}
\label{tab:wmt}
\end{table}

\subsection{FLORES-200 Multilingual}

\begin{table}[H]
\centering
\caption{FLORES-200 average spBLEU by resource level}
\begin{tabular}{lccc}
\toprule
\textbf{Resource Level} & \textbf{Language Count} & \textbf{Zen-Translator} & \textbf{Direct-Only Baseline} \\
\midrule
High-resource ($>$10M pairs)  & 42 & 46.8 & 44.1 \\
Mid-resource (1M--10M)        & 68 & 38.4 & 34.7 \\
Low-resource (100K--1M)       & 52 & 29.1 & 21.3 \\
Very low-resource ($<$100K)   & 38 & 18.7 & 8.4  \\
\textbf{Overall average}      & 200 & \textbf{38.4} & \textbf{29.6} \\
\bottomrule
\end{tabular}
\label{tab:flores}
\end{table}

\subsection{Document-Level Coherence}

\begin{table}[H]
\centering
\caption{Document-level coherence evaluation (vs. sentence-level baseline)}
\begin{tabular}{lcc}
\toprule
\textbf{Metric} & \textbf{Zen-Translator (doc)} & \textbf{Zen-Translator (sent)} \\
\midrule
Pronoun coreference accuracy (\%)  & 89.4 & 71.2 \\
Discourse connective recall (\%)   & 84.7 & 63.8 \\
Terminology consistency (\%)       & 93.1 & 78.4 \\
Overall document COMET             & 0.871 & 0.842 \\
Human preference (doc coherence)   & 4.1/5 & 3.2/5 \\
\bottomrule
\end{tabular}
\label{tab:doc}
\end{table}

\subsection{Domain Adaptation}

\begin{table}[H]
\centering
\caption{Domain-specific translation quality (BLEU, en $\to$ de)}
\begin{tabular}{lccc}
\toprule
\textbf{Domain} & \textbf{Zen-Translator} & \textbf{General MT} & \textbf{Domain Fine-tuned} \\
\midrule
Legal       & 42.3 & 36.1 & 43.8 \\
Medical     & 40.1 & 33.4 & 41.7 \\
Financial   & 41.8 & 35.2 & 42.9 \\
Technical   & 39.4 & 32.8 & 40.6 \\
Scientific  & 38.7 & 31.9 & 39.8 \\
\textbf{Average} & \textbf{40.5} & \textbf{33.9} & \textbf{41.8} \\
\bottomrule
\end{tabular}
\label{tab:domain}
\end{table}

\subsection{Register Control}

\begin{table}[H]
\centering
\caption{Register accuracy: \% of translations rated correct register by native speakers}
\begin{tabular}{lcc}
\toprule
\textbf{Register Target} & \textbf{Register Accuracy (\%)} & \textbf{BLEU Delta vs. Neutral} \\
\midrule
Formal     & 91.2 & $-0.8$ \\
Neutral    & 94.7 & baseline \\
Informal   & 88.4 & $-1.4$ \\
\bottomrule
\end{tabular}
\label{tab:register}
\end{table}

%% -----------------------------------------------------------------------
\section{Related Work}
\label{sec:related}

\subsection{Neural Machine Translation}

Transformer-based NMT~\cite{vaswani2017attention} replaced RNN-based models and
quickly became the dominant architecture due to parallelizable training and strong
performance on high-resource language pairs. Multilingual NMT~\cite{aharoni2019massively}
demonstrated that a single model can translate between many language pairs, with
transfer learning benefits for low-resource directions.

\subsection{Document-Level Translation}

Early document MT work used extra context sentences concatenated to the input.
More principled approaches using hierarchical attention~\cite{docmt} and discourse
structure modeling improved coherence, though most production systems still operate
sentence-by-sentence. Zen-Translator's 8192-token document encoder represents one
of the largest context windows applied to production NMT.

\subsection{Low-Resource Translation}

Transfer learning from high-resource languages~\cite{zoph2016transfer} and multilingual
pivoting have been the primary tools for low-resource MT. Unsupervised MT using
backtranslation~\cite{lample2018phrase} extended coverage to language pairs with
no parallel data. Zen-Translator synthesizes these approaches in a unified model.

\subsection{Domain Adaptation}

Domain-specific fine-tuning has been the standard approach for specialized domains,
requiring separate model checkpoints per domain. Adapter-based domain specialization~\cite{adapters}
enables domain switching within a single model, analogous to Zen-Translator's
domain-aware expert routing.

%% -----------------------------------------------------------------------
\section{Conclusion}
\label{sec:conclusion}

Zen-Translator establishes a new standard for multilingual neural machine translation
through the combination of document-level context (8192-token), automatic domain adaptation,
low-resource pivot routing, and register control in a single 7B parameter model.

The FLORES-200 results are particularly notable: Zen-Translator achieves spBLEU 18.7
on very-low-resource language pairs ($<$100K training pairs), more than doubling the
baseline system's performance through pivot routing and cross-lingual transfer. This
capability has direct humanitarian value for the 3+ billion people whose primary
languages lack viable translation tools.

Future work will extend document context to 128K tokens using ring attention, add
real-time simultaneous translation (speech-in, text-out), and develop evaluation
frameworks for document coherence in low-resource languages where human evaluation
is constrained by annotator availability.

\section*{Acknowledgments}

The authors thank Zoo Labs Foundation for low-resource language annotation support
and the DSO network for distributed translation quality evaluation across 200+ languages.

\bibliographystyle{plain}
\begin{thebibliography}{9}

\bibitem{wu2016google}
Y. Wu et al.,
``Google's Neural Machine Translation System: Bridging the Gap between Human and Machine Translation,''
\textit{arXiv:1609.08144}, 2016.

\bibitem{vaswani2017attention}
A. Vaswani et al.,
``Attention Is All You Need,''
\textit{NeurIPS}, 2017.

\bibitem{aharoni2019massively}
R. Aharoni, M. Johnson, O. Firat,
``Massively Multilingual Neural Machine Translation,''
\textit{NAACL}, 2019.

\bibitem{docmt}
J. Tiedemann, Y. Scherrer,
``Neural Machine Translation with Extended Context,''
\textit{DiscoMT Workshop}, 2017.

\bibitem{zoph2016transfer}
B. Zoph et al.,
``Transfer Learning for Low-Resource Neural Machine Translation,''
\textit{EMNLP}, 2016.

\bibitem{lample2018phrase}
G. Lample et al.,
``Phrase-Based and Neural Unsupervised Machine Translation,''
\textit{EMNLP}, 2018.

\bibitem{adapters}
N. Houlsby et al.,
``Parameter-Efficient Transfer Learning for NLP,''
\textit{ICML}, 2019.

\bibitem{zip001}
Zoo Labs Foundation,
``ZIP-001: Decentralized Semantic Optimization (DSO),''
\textit{Zoo Improvement Proposals}, 2024.

\end{thebibliography}

\end{document}
